\documentclass[11pt]{article}
\usepackage[margin=1in]{geometry}
\usepackage{amsmath,amssymb,amsthm,mathtools}
\usepackage[hidelinks]{hyperref}
\newtheorem{theorem}{Theorem}
\newtheorem{lemma}[theorem]{Lemma}
\newtheorem{corollary}[theorem]{Corollary}
\newcommand{\Q}{\mathbb Q}
\newcommand{\E}{\mathbb E}
\newcommand{\C}{\mathbb C}
\newcommand{\F}{\mathbb F}
\newcommand{\Z}{\mathbb Z}
\newcommand{\vp}{v_p}
\title{Mixed-prime moment windows: exponent lifting and local square-divisor degrees}
\author{Internal proof and exact certificates; no priority claim}
\date{16 November 2026}
\begin{document}
\maketitle
\noindent\textbf{Progress and boundary.}
The \href{PRIME_POWER_TWO_TERM_SFC.tex}{prime-power theorem} is revalidated.
The first uncovered exponent, six, does not have squarefree first Newton
residuals at both primes two and three for every moment index. Here we
retain that obstruction, prove an exact prime-power insertion identity for
arbitrary cofactors, and combine residual multiplicities at different primes.
The outcome includes unbounded families of mixed exponents and moment
indices, as well as stronger finite-input certificates. It does not settle
every starting index for exponent six or unrestricted SFC.

Write
\[
 f_{N,d,\alpha}(z)=\sum_{j=0}^N\binom Nj(\alpha)_{dj}z^j,
 \qquad f_{N,d}=f_{N,d,1}.
\]
These are the actual $\operatorname{Gamma}(\alpha,1)$ moments of
$(1+zY^d)^N$. The identity
$Nf_{N-1,d,\alpha}=Nf_{N,d,\alpha}-z f'_{N,d,\alpha}$ identifies a
common adjacent root with a multiple root of the higher polynomial.
The literal Newton coordinates and residual convention are those in the
preceding proof, with coordinate $j$, not reversed $N-j$.

The single-edge input for $N=kp^r$, $mk<p$, appears in
\href{sources/filaseta-rocks-2016-published.pdf}{Filaseta--Rocks,
journal Online First copy, Lemma 2.3}, and
\href{sources/rocks-2015-author-preprint.pdf}{Rocks, Proposition 3.20}.
It is rederived below, including its residual polynomial. Credit for the
factorial-unit congruence and Newton method remains with those sources;
no novelty claim is made. The proofs and controls here have not had
outside specialist review or formal verification.

\section{Exact SFC reduction and the real-place restriction}
For $d\geq1$ an integer and a fixed positive real shape $\alpha$, write
$P_n(t)=f_{n,d,\alpha}(t)$. All roots and coefficient pairs below are
allowed to be complex. This section uses only the Gamma moment formula.

\begin{theorem}\label{thm:realparity}
The two-term SFC assertion for every nonzero $A+BY^d$ at every positive
starting index, under $\operatorname{Gamma}(\alpha,1)$, is equivalent
to squarefreeness of $P_n$ for every $n\geq1$.
For even $n$, $P_n$ has no real zero. For odd $n$, it has exactly one
real zero, which is negative and simple. Therefore every nonconstant
real polynomial $H$ with $H^2\mid P_n$ has even degree. In particular,
for rational $\alpha$, every repeated rational irreducible factor has
even degree, and $P_1,P_2,P_3$ are squarefree.
\end{theorem}
\begin{proof}
Coefficient comparison gives the exact identity
\[
 P_{n-1}(t)=P_n(t)-\frac{t}{n}P_n'(t).
\]
Since $P_n(0)=1$, a common adjacent root is nonzero, and it exists
if and only if $P_n$ has a multiple root. If $AB\ne0$, the actual
moments are $A^nP_n(B/A)$. If $A=0$ or $B=0$, all positive moments
of the nonzero phase are nonzero directly. Thus all adjacent pairs
at indices $n-1,n$, $n\geq2$, are equivalent to squarefreeness of
all $P_n$, $n\geq2$; $P_1$ is linear and already squarefree.

For real $t$, the integral
\[
 P_n(t)=\E(1+tY^d)^n,\qquad Y\sim\operatorname{Gamma}(\alpha,1),
\]
is strictly positive if $n$ is even: the integrand is nonnegative
and positive except possibly at one point of the positive half-line.
If $n$ is odd, then
\[
 P_n'(t)=n\E\big[Y^d(1+tY^d)^{n-1}\big]>0
 \qquad(t\in\mathbb R).
\]
The leading coefficient is positive, so $P_n(t)\to-\infty$ as
$t\to-\infty$, whereas $P_n(0)=1$. There is exactly one real zero,
and its positive derivative makes it simple. An odd-degree real
polynomial $H$ would have a real zero, giving a repeated real zero
of $P_n$ if $H^2\mid P_n$. This is impossible. A repeated nonreal
root of a real polynomial brings its conjugate with the same
multiplicity, requiring degree at least four. This also proves the
last assertion for degrees at most three.
\end{proof}
For $d=6$ and $\alpha=1$, $P_n$ is precisely the polynomial $F_n$
in the focused proposal. Its squarefreeness at every index is still
open here. The preceding theorem is an elementary reformulation and
real-place restriction, with no priority claim. Filaseta--Rocks,
\href{https://www.impan.pl/shop/publication/transaction/download/product/91656}
{Theorem 1.1 in the final published article}, prove lower density at
least $\log2$ for irreducibility at each fixed $d$. This is not an
all-index squarefreeness theorem.

The proposed inactive-variable endpoint limit is already proved in
\href{COUPLED_BLOCK_PERIODS.tex}{the canonical block-limit note},
Theorem 1 and its single-top-monomial corollary. That proof includes
arbitrary complex lower terms, all positive Gamma shapes, and the
uniform inactive-boundary and unbounded-tail estimates. It is not an
outstanding assumption of the present squarefreeness work.

\section{Inserting a prime power preserves the entire first residual profile}
Fix a prime $p$, arbitrary positive integers $m,a,N$, and put
\[
 d=mp^a,\qquad b=m(p^a-1)/(p-1),\qquad\gamma=(-1)^b\in\F_p^\times.
\]
Factorial splitting and Wilson give, for every $j\ge0$,
\[
 U_j=\frac{(mp^aj)!}{p^{bj}(mj)!}\in\Z_p^\times,
 \qquad U_j\equiv\gamma^j\pmod p.                 \tag{1}
\]
If $1\le\alpha\le p^a$ is an integer, then also
\[
 V_j=\frac{(\alpha)_{dj}}{(dj)!}
     =\prod_{\ell=1}^{\alpha-1}(1+dj/\ell)
     \in\Z_p^\times,\qquad V_j\equiv1\pmod p.     \tag{2}
\]
Thus the coefficient of $X^j$ in $f_{N,d,\alpha}(X/p^b)$ is
the coefficient of $X^j$ in $f_{N,m}(X)$ times $U_jV_j$.

\begin{theorem}\label{thm:lift}
These two polynomials have identical Newton polygons at $p$.
If an edge has reduced slope $h/e$ and residual $R(T)$ for the base
polynomial $f_{N,m}$, its residual for the scaled lifted polynomial is
$R(\gamma^eT)$. In particular every residual irreducible degree and
multiplicity, not just the edge slopes, is preserved.
This holds for all $N,m,a$, with no finite-index extrapolation.
\end{theorem}
\begin{proof}
Every multiplier is a unit, so valuations agree. Along edge indices
$L+je$ the normalized multiplier ratio is
$U_{L+je}V_{L+je}/(U_LV_L)\equiv\gamma^{je}$.
This is precisely scaling the residual variable by $\gamma^e$.
Nonzero variable scaling preserves factor degrees and multiplicities
over the residue field. This proves every assertion.
\end{proof}

One must not replace the hypothesis ``base residuals are squarefree'' by
``the base polynomial is squarefree in characteristic zero''. Repeated
residual factors can require higher Newton data; first-order preservation
does not preserve that uncomputed data. For example $f_{7,2}$ has at $p=3$
an edge of reduced slope one with residual $1+2T^3=(1-T)^3$,
although the rational polynomial itself is squarefree.

\section{One finite residual kernel proves an unbounded indexed family}
Fix positive $m,k$ and a prime $p>mk$. Let $N=kp^r$, $r\ge1$, and
$e=p^r$, $c=(p^r-1)/(p-1)$.

\begin{lemma}\label{lem:kernel}
The Newton polygon of $f_{N,m}$ at $p$ is a single edge from
$(0,0)$ to $(kp^r,mkc)$, with reduced slope $mc/e$. Its normalized
residual is
\[
 R(T)=\overline{f_{k,m}}((-1)^{mc}T)\in\F_p[T].   \tag{3}
\]
\end{lemma}
\begin{proof}
For $0<j<N$, set $t=\vp(j)\le r$.
As $N=kp^r$ with $k<p$, the base-$p$ subtraction rule gives
$\vp\binom Nj=r-t$; and $mj<mkp^r<p^{r+1}$.
If $t<r$, then
\[
 \vp\{\tbinom Nj(mj)!\}
 =r-t+\sum_{u=1}^r\lfloor mj/p^u\rfloor
 >\sum_{u=1}^r mj/p^u=mjc/e.
\]
For $u\le t$ the quotients are integers; each remaining floor is
strictly greater than its argument minus one, proving the strict inequality.
For $j=\ell e$, all quotients are integral and equality holds.
The endpoints give the claimed height, and $\gcd(mc,e)=1$ because
$m<p$ and $c\equiv1\pmod p$. Thus the edge lattice indices are
exactly $j=\ell e$, $0\le\ell\le k$.

The identity $(1+X)^{kp^r}=(1+X^{p^r})^k$ in characteristic $p$ gives
$\binom{kp^r}{\ell p^r}\equiv\binom k\ell$.
The factorial-unit digit formula gives
\[
 p^{-m\ell c}(m\ell p^r)!
 \equiv(-1)^{m\ell c}(m\ell)!\pmod p,
\]
since $m\ell<p$. These are exactly the coefficients in (3).
\end{proof}

\begin{theorem}\label{thm:indexed}
If $p>mk$ and $\overline{f_{k,m}}$ is squarefree over $\F_p$, then for
EVERY $a,r\ge1$ and integer $1\le\alpha\le p^a$, the two moments
at indices $N-1,N$, $N=kp^r$, of every nonzero
$A+BY^{mp^a}$ cannot both vanish under $\operatorname{Gamma}(\alpha,1)$.
The coefficients $A,B$ may be arbitrary complex numbers.
\end{theorem}
\begin{proof}
Lemma~\ref{lem:kernel} has a squarefree residual independently of $r$.
Theorem~\ref{thm:lift} preserves that property independently of $a$ and
the stated shape. The weighted-initial-product lemma in the preceding
proof excludes every repeated rational factor, hence every multiple
complex root of the lifted parameter polynomial. The adjacent derivative
identity proves the two-moment statement. The coefficient-zero cases are
direct nonzero Gamma moments or nonzero constants.
\end{proof}

For $k=1$ the kernel is $1+m!T$, automatically squarefree when $p>m$.
Thus the theorem applies for every $m<p$, without a search.
Examples with mixed prime factors include $d=2\cdot3^a$ at terminal
indices $3^r$, and $d=2\cdot7^a$ at terminal indices $2\cdot7^r$.
For the second example, $k=2$ and the kernel is
$1+4T+24T^2\equiv1+4T+3T^2\pmod7$, with nonzero discriminant.
Both $a$ and $r$ range without bound. These are indexed SFC windows,
not an assertion at every positive starting index for these phases.

The finite kernel condition is necessary for this certificate: with
$(m,k,p)=(2,2,5)$ the same kernel reduces to
$1+4T+4T^2=(1+2T)^2$. The test returns no decision there.
No zero moment or non-squarefree characteristic-zero polynomial follows.

\section{Combining square-divisor degree restrictions across primes}
Consider any $P\in\Q[X]$ with nonzero constant and leading coefficient.
For a prime $p$, factor each normalized edge residual as
\[
 R_e(T)=u_e\prod_\phi\phi(T)^{m_{e,\phi}},
\]
where $\phi$ is monic irreducible over $\F_p$ and $e$ in the formula
below denotes the denominator of that edge's reduced slope. Define
the finite set
\[
 \mathcal D_p(P)=
 \left\{\sum_{\text{edges},\phi} e\deg(\phi)\,t_{e,\phi}:
              0\le t_{e,\phi}\le\lfloor m_{e,\phi}/2\rfloor\right\}. \tag{4}
\]

\begin{theorem}\label{thm:degree}
If $H^2$ divides $P$ over $\Q_p$, then $\deg H\in\mathcal D_p(P)$.
Consequently, for any finite set of primes $S$,
\[
 \bigcap_{p\in S}\mathcal D_p(P)=\{0\}
 \quad\Longrightarrow\quad P\text{ is squarefree over }\Q\text{ and }\C.
\]
The intersection condition is sufficient, not necessary.
\end{theorem}
\begin{proof}
At any slope the residual of $H$ divides the residual of $P$ with twice
its multiplicities, by the weighted initial-product argument over $\Q_p$.
Thus its multiplicity in $\phi$ is at most the corresponding floor in (4).
Its edge width is $e$ times its residual degree. The sum of all the
edge widths of $H$ is $\deg H$, since $H(0)\ne0$. This gives (4).

If a rational polynomial is not squarefree, some nonconstant irreducible
factor in $\Q[X]$ occurs squared. That same square divisor over each
$\Q_p$ has one fixed positive degree in every $\mathcal D_p(P)$.
An intersection equal to $\{0\}$ is therefore impossible.
No comparison of the fields $\Q_p$ themselves is used.
\end{proof}

For the actual moment polynomials with rational positive shape,
Theorem~\ref{thm:realparity} sharpens the sufficient condition to
\[
 \left(\bigcap_{p\in S}\mathcal D_p(P_n)\right)
 \cap\{2,4,\ldots,2\lfloor n/4\rfloor\}=\varnothing. \tag{4a}
\]
Any repeated rational irreducible factor would have an even degree in
that range and would survive every local profile. The generic local
sets $\mathcal D_p(P)$ for arbitrary rational polynomials are unchanged;
the real-place filter applies specifically to the proved moment family.
Thus a local residual may have repeated factors yet the moment
polynomial can be certified squarefree using just that prime and (4a).
An all-index exponent-six proof by this route must establish (4a) for
\emph{every} $n\geq2$, using suitable primes. A bounded table or an
irreducibility-density statement does not prove this coverage.

For $P=f_{7,6}$, the exact profiles give
$\mathcal D_2(P)=\{0,2\}$ and $\mathcal D_3(P)=\{0,1\}$.
Neither prime alone has all residuals squarefree; their local degree
intersection certifies the moment pair at indices six and seven.
With the real-place filter, prime three alone suffices because its
only possible positive square-divisor degree is one.
For $P=f_{22,6}$ the sets are $\{0,8\}$ and
$\{0,1,2,3,4\}$, again giving a certificate only after combination.
By contrast, at terminal index 21 the two primes leave
$\{0,2\}$ after intersection. Both terminal indices 21 and 23 retain degree
two, so the real-place filter leaves them unresolved as well. That is
an unresolved local test, not a counterexample. A third prime, five, certifies that particular index.
These literal examples are finite certificates; they prove no blanket
bound on the next useful prime or closure of all exponent-six indices.

The exporter also accepts any positive rational Gamma shape for a supplied
finite index and exponent, using its actual rational moment coefficients.
The generic degree theorem applies unchanged. This is not a universal
all-index Gamma-shape theorem.

\section{Ordinary pair ideals and original-marker consequences}
Every certified adjacent pair at terminal $N\ge2$, whether supplied by
the indexed family or the local degree intersection, has no common
projective coefficient zero. Thus its homogeneous moment ideal in
$\Q[A,B]$ is a regular sequence of degrees $N-1,N$. As in the preceding
proof its ordinary quotient has
\[
 \text{length }N(N-1),\qquad
 \text{top degree }2N-3\text{ of dimension one}.
\]
Over EVERY commutative $\Q$-algebra two zero moments at those indices
force $(A,B)^{2N-2}=0$. This bound is sharp in the universal quotient,
and the polynomial phase's $(2N-2)$nd power is zero. Every original
polynomial marker then has zero moments at all later exponents. The
same ordinary pair data transfers to the even-block Gaussian lift when
the Gamma shape is integral. This is nilpotence over rings, not
coefficientwise equality to zero there.

\section{Exact controls and surviving task}
\href{mixed_prime_windows.py}{The implementation} verifies literal
insertion identities, recognizes the indexed family by a finite residual
gcd, and exports multiple-prime degree profiles for the actual moment
polynomial, retaining both raw local and real-place-filtered degrees. Residual factors are checked by reconstructing the polynomial
and by the finite-field irreducibility criterion using Frobenius remainders,
separately from the factorization routine. Outside tests return unresolved.

Controls compare literal large polynomials with the small-kernel formulas,
separate rational gcd/B\'ezout computations, Gamma moments, mixed-prime
examples, repeated-factor countercontrols and ordinary quotient profiles.
All-unbounded-$a,r$ conclusions use the written formulas, not the controls.
Full exponent-six SFC, unrestricted two-term and general SFC, unrestricted
FC in at least three variables, arbitrary coupled angular phases and the
remaining source audits are still open. The inactive-variable extension is closed; the all-index squarefreeness
and broader SFC problems remain open.
\end{document}
